\documentclass[11pt]{article}
\usepackage[margin=25mm]{geometry}
\usepackage{fontspec,unicode-math}
\setmainfont{Latin Modern Roman}
\setmathfont{Latin Modern Math}
\usepackage{amsmath,amssymb,mathtools}
\usepackage{xurl,hyperref}
\hypersetup{colorlinks=true,urlcolor=blue,linkcolor=blue}
\setcounter{secnumdepth}{0}
\setlength{\parindent}{0pt}
\setlength{\parskip}{5pt}
\setlength{\emergencystretch}{3em}
\newcommand{\sref}[2]{\hyperlink{source:#1:#2}{[S#2]}}
\newcommand{\cataloguescope}[1]{\paragraph{Source question.}#1}
\newcommand{\floor}[1]{\left\lfloor#1\right\rfloor}
\begin{document}
\section{Hidden-Layer Depth Complexity of Piecewise-Linear Functions}\label{statement}

\subsection{Definitions and mathematical statement}

\paragraph{Shared notation.}Standard number systems are denoted by $\mathbb N,\mathbb Z,\mathbb Q,\mathbb R,\mathbb C$; index ranges, computational models and specialized notation follow the cited problem source and the definitions below.


ReLU is $\sigma(t)=\max\{0,t\}$. Hidden layers are the intermediate affine-plus-activation stages, excluding the final affine output; a continuous piecewise-linear function has finitely many affine regions with continuous agreement across boundaries.

A feed-forward neural network composes affine maps and specified nonlinear activations. ReLU is $\sigma(t)=\max\{0,t\}$. Empirical risk averages a loss on training samples; population risk is its expectation on the target distribution, and generalization compares these risks.

A continuous piecewise-linear function is continuous and agrees with an affine function on each member of a finite polyhedral partition. A difference-of-convex decomposition writes it as $g-h$ with both $g$ and $h$ convex.
\paragraph{Statement recorded in the supplied source.}
What is the exact minimum hidden-layer depth required for a ReLU network to represent the function $(x_1,\ldots,x_d)\mapsto\max\{x_1,\ldots,x_d\}$? More generally, does there exist a continuous piecewise-linear function that cannot be represented by any ReLU network with at most two hidden layers? \sref{statement}{2}
\paragraph{Scope qualification.}
Width, weight field, and layer-count convention are those of the source. The two questions are both retained, without equating a maximum-specific lower bound with a general depth separation.
\paragraph{Source status.}
The supplied catalogue labels this record open. Its source-status metadata was checked against the supplied snapshot on 2026-10-11. This import does not independently establish current open status or mathematical correctness. The catalogue literature assessment is dated 2026-08-22.
\subsection{Short English statement}
Find the exact depth needed to represent maximum and determine whether any piecewise-linear function needs more than two hidden layers.
\subsection{Sources}
\begin{itemize}
\item[\hypertarget{source:statement:1}{S1}] Ulam AI and the UnsolvedMath Contributors, UnsolvedMath: A Curated Collection of Open Mathematics Problems, catalogue entry. Catalogue curation and metadata: CC BY 4.0. The supplied source question is retained; definitions and the English summary are editorial additions.
\url{https://huggingface.co/datasets/ulamai/UnsolvedMath}.
\item[\hypertarget{source:statement:2}{S2}] Modern and Emerging Phenomena in Machine Learning (2026), Oberwolfach Reports, DOI 10.4171/owr/2026/11, pp. 683--736
\url{https://doi.org/10.4171/owr/2026/11}.
\item[\hypertarget{source:statement:3}{S3}] Ruess et al., Shallower ReLU Network Representations
\url{https://arxiv.org/abs/2607.21651}.
\item[\hypertarget{source:statement:4}{S4}] Bakaev et al., Better Neural Network Expressivity
\url{https://arxiv.org/abs/2505.14338}.
\textit{Underlying Oberwolfach Reports material retains its source-specific licence, including CC BY-SA 4.0 where applicable; see THIRD\_PARTY\_NOTICES.md.}
\end{itemize}
\paragraph{Preparation and review.}
Original question attributed to the cited source and UnsolvedMath contributors. Codex (GPT-6; exact serving revision not exposed) assisted with definitions, English summaries, source inspection and formatting. No new human mathematical authorship or independent proof verification is asserted. Catalogue identity checked against the complete supplied registry and existing repository statements; known flagged repetitions and obvious identical conjectures excluded. Full mathematical-equivalence/literature review has not been performed. Statement preparation and final review followed contribution prompts C and D; automated checks verify structure and identity, not open status.
\label{end:statement}
\end{document}
